\section{Negasi dan Hukum De Morgan Berkuantor}

Negasi pada kalimat berkuantor mengikuti aturan saling tukar kuantor \cite{rosen2019}. Sebagaimana pada logika proposisi, prinsip De Morgan juga berlaku pada logika predikat dengan bentuk yang disesuaikan untuk kuantor.

\subsection{Penyangkalan (Negasi) $\forall$}
Negasi dari pernyataan universal \textit{``Semua mahasiswa di kelas ini pandai''} bukanlah \textit{``Semua mahasiswa tidak pandai''}, melainkan \textit{``Ada mahasiswa yang tidak pandai''}.

Secara formal:
\[ \neg \forall x \, P(x) \equiv \exists x \, \neg P(x) \]

\subsection{Penyangkalan (Negasi) $\exists$}
Negasi dari pernyataan eksistensial \textit{``Ada dosen di kampus ini yang tegas berlebihan''} dapat dinyatakan sebagai \textit{``Tidak ada dosen yang tegas berlebihan''}, yang ekuivalen dengan \textit{``Semua dosen tidak tegas berlebihan''}.

Secara formal:
\[ \neg \exists x \, P(x) \equiv \forall x \, \neg P(x) \]

\subsection{Aplikasi Matematis Tingkat Lanjut}

Untuk negasi dengan kuantor bersarang, tanda $\neg$ didorong ke dalam secara bertahap; setiap kuantor yang dilalui berganti pasangan: $\forall \leftrightarrow \exists$.

\begin{contoh}
Negasi dari pernyataan bersarang: $\exists x \, \forall y \, \forall z \, P(x,y,z)$

\begin{align*}
    \neg \Big[ \exists x \ \forall y \ \forall z \ P(x,y,z) \Big] 
    &\equiv \forall x \ \neg \Big[ \forall y \ \forall z \ P(x,y,z) \Big] \\
    &\equiv \forall x \ \exists y \ \neg \Big[ \forall z \ P(x,y,z) \Big] \\
    &\equiv \forall x \ \exists y \ \exists z \ \neg P(x,y,z)
\end{align*}
\end{contoh}

\begin{contoh}
Untuk pernyataan kondisional ``Semua mahasiswa baru mengikuti orientasi'', translasi formalnya adalah $\forall x \, (M(x) \rightarrow O(x))$.

Gunakan ekuivalensi implikasi $p \rightarrow q \equiv \neg p \lor q$ untuk mendorong negasi.
\begin{align*}
    \neg \Big( \forall x \, (\,M(x) \rightarrow O(x)\,) \Big) 
    &\equiv \exists x \, \neg \Big( M(x) \rightarrow O(x) \Big) \\
    &\equiv \exists x \, \neg \Big( \neg M(x) \lor O(x) \Big) & \text{(Implikasi)} \\
    &\equiv \exists x \, \Big( \neg(\neg M(x)) \land \neg O(x) \Big) & \text{(De Morgan)} \\
    &\equiv \exists x \, \Big( M(x) \land \neg O(x) \Big)
\end{align*}
\textbf{Makna negasi:} ``Terdapat mahasiswa baru yang tidak mengikuti orientasi.''
\end{contoh}
